\section{Color from an 't Hooft twist}
\label{sec:colortwist}
Section~\ref{sec:gaugefieldscalars} ended with gauge groups $U(15)$ or
$U(6)\times U(9)$, in which $15=5\times3$ is a choice of basis. This
section examines one mechanism that makes the product structure
physical: a flux background of the gauge field on a two-torus, the
setting of Ref.~\cite{CDS} at rational noncommutativity. Statements
about twisted bundles and index numbers are standard and are rederived
in blackboard~013; the conclusions for the sBootstrap are conditional.

\paragraph{Twist.}
A $U(15)$ bundle of first Chern number $m=3m'$, with $m'$ prime to 5,
has transition functions $U\otimes1_3$ and $V^{m'}\otimes1_3$ on adjoint
fields, with $U,V$ the $5\times5$ clock and shift matrices. Their
commutant is $1_5\otimes U(3)$: color is the unbroken group and the
five-fold index is a wrapping number.

\paragraph{Mesons on a noncommutative torus.}
Expanding an adjoint field in $U^aV^b\otimes\phi_{ab}(x)$ gives the
$3\times3$ fields $\phi_{ab}$ fractional momenta $b/5$ and $-m'a/5$ in
units of $2\pi\hbar/L_i$. The single-valued generators
$Z_1=Ve^{2\pi ix_1/5L_1}$ and $Z_2=U^{-s}e^{2\pi ix_2/5L_2}$, with
$sm'=1$ mod 5, satisfy $Z_1Z_2=e^{2\pi is/5}Z_2Z_1$. The meson-like
fields, read in Ref.~\cite{Rivero2024} as $(24+1)\otimes(8+1)$, are
therefore $U(3)$ fields on the noncommutative torus with $\theta=s/5$:
the number of colors is the rank and the number of flavors is the
denominator of $\theta$. Flavor symmetry is reduced to
$\mathbb Z_5\times\mathbb Z_5$, and the lowest state of class $(a,b)$
has
\begin{equation}
M=\frac{2\pi\hbar}{c_{\rm light}}
\sqrt{\Bigl(\frac{\tilde b}{5L_1}\Bigr)^2+\Bigl(\frac{\tilde a}{5L_2}\Bigr)^2},
\end{equation}
with $\tilde a,\tilde b\in\{-2,\dots,2\}$ the representatives of $-m'a$
and $b$. On a square torus the 24 non-singlet classes split as
$4,4,4,8,4$ at $M^2\propto1,2,4,5,8$; only the $U(3)$ gauge multiplet
keeps a zero mode.

\paragraph{Diquark-like zero modes.}
Since $\Lambda^2(5\otimes3)=S^2(5)\otimes\Lambda^2(3)+
\Lambda^2(5)\otimes S^2(3)$, the numbers of color-antitriplet and
color-sextet zero modes on the torus are $c_1(S^2E_5)=6m'$ and
$c_1(\Lambda^2E_5)=4m'$, in place of the flavor multiplicities 15 and
10. For wrapping number $n$ they are $(n+1)m'$ and $(n-1)m'$.

\paragraph{Charge.}
For any background of the form $f\otimes1_3$ the unbroken group is the
unitary group of $S'\otimes M_3$, with $S'$ the commutant of the flavor
parts. A single color group needs $S'=\mathbb C$, which leaves no
flavor-dependent charge. A charge with eigenvalues $2/3$ and $-1/3$
needs $S'\supseteq\mathbb C\oplus\mathbb C$ and gives
$U(3)_u\times U(3)_d$. The corresponding construction has two stacks
with wrapping numbers 2 and 3, fluxes $m_u$ and $m_d$, the label count
$3\cdot2+3\cdot3+1=16$, and $Q=\tfrac23Q_u-\tfrac13Q_d+q_nQ_n$. Its
zero modes on the flux torus are $3|m_u|$ antitriplets and $|m_u|$
sextets in the $uu$ sector, $4|m_d|$ and $2|m_d|$ in the $dd$ sector,
$|3m_u+2m_d|$ $ud$ diquarks and $|3m_u-2m_d|$ charged mesons. For
$m_u=1$ the $uu$ sector has the three antitriplets of
Ref.~\cite{Rivero2024}; the counts six and six of the other sectors
would need $m_d=3/2$. A scan of factorized fluxes on $T^6$ with flux
numbers between $-4$ and $4$ finds no assignment with the ratio
$1:2:2$ and none that cancels the tadpoles with these three stacks.

\paragraph{Assessment.}
The flux mechanism singles out color by turning flavor into wrapping.
The multiplicities that Ref.~\cite{Rivero2024} counts as dimensions of
flavor representations become index numbers and momentum classes with
different values, and a flavor-dependent charge requires two color
groups. The counts of Ref.~\cite{Rivero2024} are recovered when the flux
is absent, where color is not singled out. A background outside the
flavor directions, or a mechanism other than symmetry breaking, would
be needed to obtain both.
